\documentclass[10pt,a4paper]{amsart}
\usepackage{amsmath,amssymb,amsthm,graphicx}
% Retained preamble for the Stacks Project chapter snapshot in this directory.
% The chapter's own preamble.tex is a copy of the shared Stacks Project preamble; the block
% below is sliced MECHANICALLY, line for line, from it. Nothing was added, removed or
% redefined; the wrapper class, the xr-hyper/hyperref wiring and the externaldocument list
% are omitted because the assembler supplies the document class.
%
% lines taken from preamble.tex: 3, 5, 9, 11, 12, 156, 157, 158, ... (51 lines total)
%
% package coverage restored from the chapter's own preamble (2026-10-13): xy, xy-2cell

\usepackage[all]{xy}
\xyoption{2cell}
\UseAllTwocells

\newenvironment{reference}{\comment}{\endcomment}
\newenvironment{slogan}{\comment}{\endcomment}
\newenvironment{history}{\comment}{\endcomment}
\theoremstyle{plain}
\newtheorem{theorem}[subsection]{Theorem}
\newtheorem{proposition}[subsection]{Proposition}
\newtheorem{lemma}[subsection]{Lemma}
\theoremstyle{definition}
\newtheorem{definition}[subsection]{Definition}
\newtheorem{example}[subsection]{Example}
\newtheorem{exercise}[subsection]{Exercise}
\newtheorem{situation}[subsection]{Situation}
\theoremstyle{remark}
\newtheorem{remark}[subsection]{Remark}
\newtheorem{remarks}[subsection]{Remarks}
\numberwithin{equation}{subsection}
\def\lim{\mathop{\mathrm{lim}}\nolimits}
\def\colim{\mathop{\mathrm{colim}}\nolimits}
\def\Spec{\mathop{\mathrm{Spec}}}
\def\Hom{\mathop{\mathrm{Hom}}\nolimits}
\def\Ext{\mathop{\mathrm{Ext}}\nolimits}
\def\SheafHom{\mathop{\mathcal{H}\!\mathit{om}}\nolimits}
\def\SheafExt{\mathop{\mathcal{E}\!\mathit{xt}}\nolimits}
\def\Sch{\mathit{Sch}}
\def\Mor{\mathop{\mathrm{Mor}}\nolimits}
\def\Ob{\mathop{\mathrm{Ob}}\nolimits}
\def\Sh{\mathop{\mathit{Sh}}\nolimits}
\def\NL{\mathop{N\!L}\nolimits}
\def\CH{\mathop{\mathrm{CH}}\nolimits}
\def\proetale{{pro\text{-}\acute{e}tale}}
\def\etale{{\acute{e}tale}}
\def\QCoh{\mathit{QCoh}}
\def\Ker{\mathop{\mathrm{Ker}}}
\def\Im{\mathop{\mathrm{Im}}}
\def\Coker{\mathop{\mathrm{Coker}}}
\def\Coim{\mathop{\mathrm{Coim}}}
\DeclareMathSymbol{\boxtimes}{\mathbin}{AMSa}{"02}
\def\QCohstack{\mathcal{QC}\!\mathit{oh}}
\def\Cohstack{\mathcal{C}\!\mathit{oh}}
\def\Spacesstack{\mathcal{S}\!\mathit{paces}}
\def\Quotfunctor{\mathrm{Quot}}
\def\Hilbfunctor{\mathrm{Hilb}}
\def\Curvesstack{\mathcal{C}\!\mathit{urves}}
\def\Polarizedstack{\mathcal{P}\!\mathit{olarized}}
\def\Complexesstack{\mathcal{C}\!\mathit{omplexes}}
\def\Pic{\mathop{\mathrm{Pic}}\nolimits}
\def\Picardstack{\mathcal{P}\!\mathit{ic}}
\def\Picardfunctor{\mathrm{Pic}}
\def\Deformationcategory{\mathcal{D}\!\mathit{ef}}
\usepackage{cleveref}
\crefname{originallemma}{original lemma}{original lemmas}
\Crefname{originallemma}{Original lemma}{Original lemmas}
\makeatletter
\newlabel{etale-cohomology-lemma-ext-modules-hom}{{0DVE}{source}{}{sourcepointer-1}{}}
\newlabel{etale-cohomology-lemma-ext-modules-hom@cref}{{[originallemma][0][]0DVE}{[1][0][]source}{}{}{}}
\makeatother
\crefname{originallemma}{original lemma}{original lemmas}
\Crefname{originallemma}{Original lemma}{Original lemmas}
\makeatletter
\newlabel{etale-cohomology-lemma-finite-dim-group-cohomology}{{0DVF}{source}{}{sourcepointer-2}{}}
\newlabel{etale-cohomology-lemma-finite-dim-group-cohomology@cref}{{[originallemma][0][]0DVF}{[1][0][]source}{}{}{}}
\makeatother
\crefname{originallemma}{original lemma}{original lemmas}
\Crefname{originallemma}{Original lemma}{Original lemmas}
\makeatletter
\newlabel{example-representations}{{0D3J}{source}{}{sourcepointer-3}{}}
\newlabel{example-representations@cref}{{[originallemma][0][]0D3J}{[1][0][]source}{}{}{}}
\makeatother
\crefname{originallemma}{original lemma}{original lemmas}
\Crefname{originallemma}{Original lemma}{Original lemmas}
\makeatletter
\newlabel{lemma-finite-projective-modules-TI}{{0DVQ}{source}{}{sourcepointer-4}{}}
\newlabel{lemma-finite-projective-modules-TI@cref}{{[originallemma][0][]0DVQ}{[1][0][]source}{}{}{}}
\makeatother
\title{Stacks, Examples of Deformations: lemma-representations-TI}
\author{The Stacks Project Authors}
\date{stacks-project:examples-defos.tex:local-snapshot-20261012; GNU FDL 1.2 or later}
\begin{document}
\maketitle
\noindent\textit{Editorial note.} Statement span is lines 563-578 of sources/stacks-examples-defos/pinned-original-source.tex; proof span is lines 580-623, the complete author proof (42 body lines).\par\medskip
\noindent\textit{Editorial note.} Navigation stubs: 4 cited-but-not-included labels are supplied as non-load-bearing pointers with their pinned tags where the snapshot tag index has them.\par\medskip
\noindent\textit{Editorial note.} extra\_packages is empty on purpose: the retained preamble already loads xy with options and a second declaration clashes in the assembler wrapper.\par\medskip
\noindent\textit{Editorial note.} Text is retained verbatim, including the author's own typos and placeholders.\par\medskip
\noindent\textit{Editorial note.} Queue index 14 of queue\_examples-defos.json (pb=1584, stmt 563-578, proof 580-623).\par\medskip
\noindent\textit{Allowed standard prerequisites (editorial).}\begin{itemize}
\item Explicit prerequisite boundary. The cited inputs are prerequisites; the argument in the retained span is not replaced by any cited result.
\end{itemize}\medskip
\ifdefined\hypertarget\hypertarget{sourcepointer-1}{}\fi
\noindent\textit{Source reference.} Original lemma 0DVE: sources/stacks-examples-defos/pinned-original-source.tex, source lines 623--623. Cited by the retained statement or proof but not part of the collected spans; supplied only as a navigation pointer. The argument does not depend on its content.\par\medskip
\ifdefined\hypertarget\hypertarget{sourcepointer-2}{}\fi
\noindent\textit{Source reference.} Original lemma 0DVF: sources/stacks-examples-defos/pinned-original-source.tex, source lines 623--623. Cited by the retained statement or proof but not part of the collected spans; supplied only as a navigation pointer. The argument does not depend on its content.\par\medskip
\ifdefined\hypertarget\hypertarget{sourcepointer-3}{}\fi
\noindent\textit{Source reference.} Original lemma 0D3J: sources/stacks-examples-defos/pinned-original-source.tex, source lines 483--494. Cited by the retained statement or proof but not part of the collected spans; supplied only as a navigation pointer. The argument does not depend on its content.\par\medskip
\ifdefined\hypertarget\hypertarget{sourcepointer-4}{}\fi
\noindent\textit{Source reference.} Original lemma 0DVQ: sources/stacks-examples-defos/pinned-original-source.tex, source lines 430--439. Cited by the retained statement or proof but not part of the collected spans; supplied only as a navigation pointer. The argument does not depend on its content.\par\medskip
\begin{lemma}
\label{lemma-representations-TI}
In Example \ref{example-representations} let 
$\rho_0 : \Gamma \to \text{GL}_k(V)$
be a finite dimensional representation. Then
$$
T\Deformationcategory_{V, \rho_0} = \Ext^1_{k[\Gamma]}(V, V) =
H^1(\Gamma, \text{End}_k(V))
\quad\text{and}\quad
\text{Inf}(\Deformationcategory_{V, \rho_0}) = H^0(\Gamma, \text{End}_k(V))
$$
Thus $\text{Inf}(\Deformationcategory_{V, \rho_0})$
is always finite dimensional
and $T\Deformationcategory_{V, \rho_0}$ is finite dimensional
if $\Gamma$ is finitely generated.
\end{lemma}

\begin{proof}
We first deal with the infinitesimal automorphisms.
Let $M = V \otimes_k k[\epsilon]$ with induced action
$\rho_0' : \Gamma \to \text{GL}_n(M)$.
Then an infinitesimal automorphism, i.e., an element of
$\text{Inf}(\Deformationcategory_{V, \rho_0})$,
is given by an automorphism
$\gamma = \text{id} + \epsilon \psi : M \to M$
as in the proof of Lemma \ref{lemma-finite-projective-modules-TI},
where moreover $\psi$ has to commute
with the action of $\Gamma$ (given by $\rho_0$).
Thus we see that
$$
\text{Inf}(\Deformationcategory_{V, \rho_0}) = H^0(\Gamma, \text{End}_k(V))
$$
as predicted in the lemma.

\medskip\noindent
Next, let $(k[\epsilon], M, \rho)$ be an object of $\mathcal{F}$
over $k[\epsilon]$ and let $\alpha : M \to V$ be a $\Gamma$-equivariant map
inducing an isomorphism $M/\epsilon M \to V$.
Since $M$ is free as a $k[\epsilon]$-module we obtain
an extension of $\Gamma$-modules
$$
0 \to V \to M \xrightarrow{\alpha} V \to 0
$$
We omit the detailed construction of the map on the left.
Conversely, if we have an extension of $\Gamma$-modules as
above, then we can use this to make a $k[\epsilon]$-module
structure on $M$ and get an object of $\mathcal{F}(k[\epsilon])$
together with a map $\alpha$ as above.
It follows that
$$
T\Deformationcategory_{V, \rho_0} = \Ext^1_{k[\Gamma]}(V, V)
$$
as predicted in the lemma. This is equal to
$H^1(\Gamma, \text{End}_k(V))$ by
\'Etale Cohomology, Lemma \ref{etale-cohomology-lemma-ext-modules-hom}.

\medskip\noindent
The statement on dimensions follows from
\'Etale Cohomology, Lemma
\ref{etale-cohomology-lemma-finite-dim-group-cohomology}.
\end{proof}

\end{document}
